\subsection{Competition creates competition}\label{app:prop2}

This subsection proves Proposition~\ref{prop:ccc} for general classes and priors. It then states and proves
Proposition~\ref{prop:region}, the admissible region at the benchmark. Write \(F_1\) for a strengthening of \(F_0\) and \(F_2\) for a
strengthening of \(F_1\). For \(i\in\{0,1,2\}\), write \(B_i=B_{F_i}\), \(\Delta_i=\DT(F_i)\), and \(t_{0,i}\),
\(t_{j,i}\), \(g_{j,i}\), \(w_i(\mu)\) for the objects at \(F_i\). Also write \(M=M(\pi)\), \(m=m(\pi)\) and
\(\kappa=\kappa(\pi)\). Under Assumption~\ref{ass:classes} every value meets the reserve, so
\(W_H(F)=\E_H[F(\theta)]\), the probability that a prepared high-class challenger outbids the incumbent. Value ties
have probability zero because \(F\) is continuous.

\begin{proof}[Proof of Proposition~\ref{prop:ccc}]
The proof has seven steps. Steps 1, 3 and 5 state the equilibrium tuple of each economy, check every condition
against it, and prove uniqueness from bounds on every candidate.

\emph{Step 0: implied inequalities.} Proposition~\ref{prop:returns}(a) gives \(B_1(\pi)\le B_0(\pi)\), so (H1)
gives \(B_1(\pi)<c_H<B_1(M)\). Since \(M>\pi\), \(B_1(M)-B_1(\pi)=(M-\pi)(g_{H,1}-g_{L,1})>0\). Hence
\(g_{H,1}>g_{L,1}\), the threshold \(\tau_1\) at \(F_1\) is defined, and \(\pi<\tau_1<M\). By (H2),
\(\Delta_1>k/(\rho\kappa)>0\). Proposition~\ref{prop:returns}(b) gives \(\Delta_2\ge\Delta_1\), so
\(\rho\kappa\Delta_2>k\). Proposition~\ref{prop:returns} also gives \(g_{H,2}\ge g_{L,2}\), so
\(B_2(\mu)\le B_2(M)<c_H\) for every \(\mu\le M\). The zero-cost challenger prepares after every price, so
Lemma~\ref{lem:prices}(a) gives \(e\ge\rho\) in every candidate of each economy.

\emph{Step 1: the weak incumbent.} Fix any candidate at \(F_0\). By Lemma~\ref{lem:trading}(a) and the upper bounds
in \eqref{eq:a2-bounds}, each correctly signed order earns at most \(\Delta_0-k<0\), and each wrong-signed order
earns at most \(-k\). Abstention is therefore the unique maximizer of \(U_H\) and of \(U_L\), and (E5) gives
\(\sigma_H=\sigma_L=\delta_0\) under both investor rules. Then \(f_H(x)=f_L(x)=(1-\lambda)\eta(x)+\lambda\1\{x=0\}\),
so (E1) gives \(\mu(x)=\pi\) at every order, and (E2) gives \(\bar\mu(\bar P)=\pi\) at every price. Since
\(B_0(\pi)<c_H\), (E3) gives \(\zeta(\bar P)=0\) with no tie, so \(e(x)=\rho\). Condition (E4) and
\eqref{eq:price} then fix one price for every order. Every equilibrium at \(F_0\) is therefore the tuple
\[
\mathcal E_0:\quad
\begin{aligned}
&\sigma_H=\sigma_L=\delta_0,\quad \mu(x)=\pi,\quad P(x)=\bar P_0:=t_{0,0}+\rho\,[t_{L,0}-t_{0,0}+\Delta_0\pi],\\
&\bar\mu(\bar P_0)=\pi,\quad \zeta(\bar P_0)=0,
\end{aligned}
\]
with truthful bids.

\emph{Verification of \(\mathcal E_0\).} (E1): \(f_H=f_L\), so the posterior at every order is \(\pi\). (E2): the
single price pools all orders, and the odds of \(\bar\mu(\bar P_0)\) are
\(\pi\sum_xf_H(x)/[(1-\pi)\sum_xf_L(x)]=\pi/(1-\pi)\). (E3): \(B_0(\pi)<c_H\), strictly. (E4): by
\eqref{eq:a2-terminal} with \(e\equiv\rho\), \(\E[V_T\mid X=x]=t_{0,0}+\rho\,w_0(\pi)=\bar P_0\). (E5): at
\(\mathcal E_0\), \(K_H=\rho\Delta_0(1-\pi)<k\) and \(K_L=\rho\Delta_0\pi<k\), so \(U_H(1)<0=U_H(0)\) and
\(U_L(-1)<0=U_L(0)\), and the wrong-signed payoffs are at most \(-k\). (E6): Appendix~\ref{app:payoffs}. Entry is
\(\rho\), and high-class ownership is \(\pi\rho W_H(F_0)\).

\emph{Step 2: the strong incumbent, orders in every candidate.} Fix any candidate at \(F_1\). By
\eqref{eq:a2-bounds}, \(K_H\ge\rho\kappa\Delta_1>k\) and \(K_L\ge\rho\kappa\Delta_1>k\). By
Lemma~\ref{lem:trading}(a)--(b), the purchase is the unique maximizer of class \(H\), and the sale is the unique
maximizer of class \(L\). Every equilibrium at \(F_1\) has \(\sigma_H=\delta_1\) and \(\sigma_L=\delta_{-1}\), under
both investor rules. These bounds use only the nonzero orders, so they hold for every \(\beta<1/2\).

\emph{Step 3: the strong incumbent, the tuple.} Lemma~\ref{lem:beliefs}(d) gives the posteriors \((m,\pi,M)\) after
the orders \((-1,0,1)\). Step 0 gives \(\Delta_1>0\) and \(g_{H,1}>g_{L,1}\), so Lemma~\ref{lem:prices} applies.
Since \(B_1(M)>c_H\), \(B_1(\pi)<c_H\) and \(B_1(m)<B_1(\pi)\), Lemma~\ref{lem:prices}(d) gives \(e(1)=1\) and
\(e(0)=e(-1)=\rho\), with no tie. Every equilibrium at \(F_1\) is therefore the tuple
\[
\mathcal E_1:\quad
\begin{aligned}
&\sigma_H=\delta_1,\quad \sigma_L=\delta_{-1},\quad (\mu(-1),\mu(0),\mu(1))=(m,\pi,M),\\
&P(x)=t_{0,1}+e(x)\,w_1(\mu(x)),\quad \bar\mu(P(x))=\mu(x),\quad \zeta(P(x))=\1\{x=1\},
\end{aligned}
\]
with \((e(-1),e(0),e(1))=(\rho,\rho,1)\) and truthful bids.

\emph{Verification of \(\mathcal E_1\).} (E1): the likelihoods \(f_H=(b,b_0,b+\lambda)\) and
\(f_L=(b+\lambda,b_0,b)\) give \((m,\pi,M)\). (E2): the stated entry is nondecreasing in the posterior and lies in
\([\rho,1]\), and \(\Delta_1>0\). By the paragraph on constructed tuples in Appendix~\ref{app:beliefs}, the three
prices are strictly increasing in the posterior, so \(P(-1)<P(0)<P(1)\). Each price has one order in its
preimage, so \(\bar\mu(P(x))=\mu(x)\). (E3): \(B_1(M)>c_H\), \(B_1(\pi)<c_H\) and \(B_1(m)<c_H\), all strict.
(E4): \eqref{eq:a2-terminal} with the stated entry. (E5): at \(\mathcal E_1\), \(K_H=\Delta_1(1-M)\ge\rho\kappa\Delta_1>k\)
and \(K_L=\rho\Delta_1m\ge\rho\kappa\Delta_1>k\). Hence \(U_H(1)>0=U_H(0)>U_H(-1)\) and
\(U_L(-1)>0=U_L(0)>U_L(1)\). (E6): Appendix~\ref{app:payoffs}.

\emph{Step 4: entry and ownership at \(F_1\).} Class \(H\) places the order \(1\) with probability
\(f_H(1)=b+\lambda\), through a liquidity purchase or an informed purchase. Class \(L\) places it with probability
\(f_L(1)=b\). Only the order \(1\) recruits the high-cost challenger. Hence
\[
e_H=\rho+(1-\rho)(b+\lambda),\qquad e_L=\rho+(1-\rho)b,\qquad
\Ent_1=\pi e_H+(1-\pi)e_L=\rho+(1-\rho)(b+\pi\lambda)>\rho ,
\]
since \(b>0\). Preparation depends on the order and the cost, and both are independent of \((\theta,R)\) given the
class. A prepared high-class challenger acquires the target when \(\theta>R\), with probability \(W_H(F_1)\). So
\(\Own_{H,1}=\pi e_HW_H(F_1)\), and Step 1 gives \(\Own_{H,0}=\pi\rho W_H(F_0)\). Division by \(\pi>0\) shows that
ownership rises exactly when \([\rho+(1-\rho)(b+\lambda)]\,W_H(F_1)>\rho\,W_H(F_0)\). In the point-mass benchmark
\(h>r\) at every strength, so \(W_H=1\), and the condition reads \(e_H>\rho\), which holds.

\emph{Step 5: the very strong incumbent.} Fix any candidate at \(F_2\). By \eqref{eq:a2-bounds},
\(K_H\ge\rho\kappa\Delta_2>k\) and \(K_L\ge\rho\kappa\Delta_2>k\), so every equilibrium has full orders, and
Lemma~\ref{lem:beliefs}(d) gives the posteriors \((m,\pi,M)\). By Step 0, \(B_2(\mu)\le B_2(M)<c_H\) at each
posterior, and every price belief is one of these posteriors or an average of them. So (E3) gives
\(\zeta\equiv0\) with no tie, and \(e\equiv\rho\). Every equilibrium at \(F_2\) is therefore the tuple
\[
\mathcal E_2:\quad
\begin{aligned}
&\sigma_H=\delta_1,\quad \sigma_L=\delta_{-1},\quad (\mu(-1),\mu(0),\mu(1))=(m,\pi,M),\\
&P(x)=t_{0,2}+\rho\,w_2(\mu(x)),\quad \bar\mu(P(x))=\mu(x),\quad \zeta\equiv0,
\end{aligned}
\]
with truthful bids. \emph{Verification.} (E1): the likelihoods \(f_H=(b,b_0,b+\lambda)\) and \(f_L=(b+\lambda,b_0,b)\)
give \((m,\pi,M)\). (E2): since \(\Delta_2>0\) and \(m<\pi<M\), the three
prices \(t_{0,2}+\rho\,w_2(\mu(x))\) differ, so each price has one order in its preimage. (E3): \(B_2(\mu(x))<c_H\)
at each posterior. (E4): \eqref{eq:a2-terminal} with \(e\equiv\rho\). (E5): at \(\mathcal E_2\),
\(K_H=\rho\Delta_2(1-M)>k\) and \(K_L=\rho\Delta_2m>k\). (E6): Appendix~\ref{app:payoffs}. The prices differ, so the
price stays informative while entry equals \(\rho\).

\emph{Step 6: tie rules.} Every investor comparison in Steps 1, 2 and 5 is strict. Every high-cost comparison at an
equilibrium belief is strict. No tie rule is used, so the four pairs of rules give the same unique outcome in each
economy.
\end{proof}

\paragraph{Each hypothesis is used once.}
\begin{itemize}[leftmargin=2em]
\item \(B_0(\pi)<c_H\) gives entry \(\rho\) at \(F_0\). Through Proposition~\ref{prop:returns}(a), it also stops
recruitment after the zero order at \(F_1\).
\item \(c_H<B_1(M)\) gives recruitment after a purchase at \(F_1\).
\item \(\Delta_0<k\) defeats every informed order at \(F_0\), against every candidate.
\item \(k<\rho\kappa\Delta_1\) forces full orders at \(F_1\), and through Proposition~\ref{prop:returns}(b) at \(F_2\),
against every candidate.
\item \(B_2(M)<c_H\) stops recruitment at \(F_2\).
\end{itemize}
The bound at the zero order, Lemma~\ref{lem:beliefs}(c), is the only use of \(\beta\le1/3\), and it enters no step.

The weaker condition \(B_1(\pi)<c_H\) cannot replace \(B_0(\pi)<c_H\). Keep calibration B and the benchmark values,
but set \(c_H=\reg{nec_cH}\) and \(r_2=\reg{nec_r2}\). Then \(k-\Delta_0=\reg{nec_margin_kweak}\),
\(c_H-B_1(1/2)=\reg{nec_margin_prior_r1}\), \(B_1(M)-c_H=\reg{nec_margin_recruit_B}\),
\(\rho\kappa\Delta_1-k=\reg{nec_margin_kstrong_B}\) and \(c_H-B_2(M)=\reg{nec_margin_return_B}\) are positive, but
\(c_H-B_0(1/2)=\reg{nec_margin_weakcost}\). At \(r_0\), Step 1 gives zero orders and the posterior \(1/2\), and
\(B_0(1/2)=\reg{B_r0_half}>c_H\). Both cost types prepare, so entry is \(\reg{nec_entry_r0}\) at \(r_0\). Entry then
falls to \(\reg{entry_r1_B}\) at \(r_1\). On the benchmark path the left side of (H2) also guards a conclusion. At
\(r=\reg{str_r155}\) it fails, since \(\DT=\reg{Delta_r155_frac}>k\) there, and Proposition~\ref{prop:coexistence}
gives two equilibria.

\paragraph{The tuples at calibration B.}
Table~\ref{tab:outcomes} lists \(\mathcal E_0\), \(\mathcal E_1\) and \(\mathcal E_2\) at calibration B and the prior
\(\pi=\reg{prim_pi}\). In \(\mathcal E_1\) the posteriors are \((\reg{m_half_B},\,\pi,\,\reg{M_half_B})\) and the
prices are \(\reg{price_r1_m1_B}\), \(\reg{price_r1_0_B}\) and \(\reg{price_r1_p1_B}\). In \(\mathcal E_2\) the prices
are \(\reg{price_r2_m1_B}\), \(\reg{price_r2_0_B}\) and \(\reg{price_r2_p1_B}\), and in \(\mathcal E_0\) every order
has the price \(\reg{price_r0}\). The values come from
\nolinkurl{P1_core/p1_numbers.py} (run log
\nolinkurl{P1_core/p1_numbers.log}). This script evaluates every quantity
that depends on \(\eu\) in outward interval arithmetic (\texttt{mpmath.iv}, 50 digits). The independent
recomputation \nolinkurl{XP1/xp1_numbers.py} (run log
\nolinkurl{XP1/xp1_numbers.log}) bounds \(\eu\) by its series through the
term \(1/110!\) with a rational bound on the tail. It records exact rational endpoints for every printed value in
\nolinkurl{XP1/xp1_registry_recompute.csv}. These rational endpoints are
the certificate of each display.

\paragraph{The open set.}
For uniform incumbents, (H1) to (H3) are strict inequalities between continuous functions of
\((r_0,r_1,r_2,c_H,k,\rho,\pi,\lambda,\beta)\). They define an open set in the domain \(0<\beta<1/2\). The model
adds the structural bound \(\beta\le1/3\), and in that domain the set is relatively open.
Proposition~\ref{prop:region} gives its section through the benchmark. The sufficient condition
\((1+\rho)W_H(F_1)\ge2\rho W_H(F_0)>0\) for higher ownership implies the exact condition of Step 4 whenever
\(b+\lambda>1/2\), because then \(e_H>(1+\rho)/2\). At calibration B, \(b+\lambda=\reg{b_plus_lambda_B}\).

\begin{proposition}[The admissible region; analytical]\label{prop:region}
Take the benchmark auction and cost primitives, the strengths \(r_0,r_1,r_2\), and a slot with \(0<\lambda<1\) and
\(0<\beta\le1/3\).
\begin{enumerate}[label=(\alph*),leftmargin=2em]
\item At the prior \(1/2\), (H1)--(H3) hold if and only if
\[
\reg{u_lo_frac}<u<\reg{u_hi_frac},\quad\text{that is,}\quad
\tau(r_1)=\reg{M_lo_frac}<M(1/2)<\reg{M_hi_frac}=\tau(r_2).
\]
Since \(\lambda=u\beta/(1+u\beta)\) rises with \(u\), the band is an open interval of \(\lambda\) at each
\(\beta\).
\item The open band is exactly the set on which the strict hypotheses hold. If \(u<\reg{u_lo_frac}\), entry at
\(r_1\) is \(\rho\). At \(u=\reg{u_lo_frac}\), orders are uniquely zero at \(r_0\) and uniquely full at \(r_1\) and
\(r_2\). Rule C1 gives every conclusion of Proposition~\ref{prop:ccc}. Under rule CM, entry after a purchase at \(r_1\) can take any value
in \([\rho,1]\). If \(\reg{u_hi_frac}\le u<\reg{u_kstrong_frac}\), entry after a purchase at \(r_2\) is one under
rule C1. If \(u\ge\reg{u_kstrong_frac}\), the right side of (H2) fails at \(r_1\).
\item With \(1+u=\eu\) and the other benchmark inputs fixed, (H1)--(H3) hold if and only if
\[
\frac{\tau(r_1)}{\tau(r_1)+[1-\tau(r_1)]\,\eu}<\pi<\frac{\tau(r_2)}{\tau(r_2)+[1-\tau(r_2)]\,\eu},
\]
an interval with endpoints \(\approx\reg{prior_lo_B}\) and \(\approx\reg{prior_hi_B}\). Hypotheses (H1) and (H2)
give parts (i) and (ii) of Proposition~\ref{prop:ccc}. They hold if and only if \(\pi\) lies above the same lower
endpoint and below \(\tau(r_0)=\reg{tau_r0_frac}\approx\reg{prior_ii_hi}\).
\end{enumerate}
\end{proposition}

\begin{proof}[Proof of Proposition~\ref{prop:region}]
Take the benchmark primitives \((h,\ell,p)=(\reg{prim_h},\reg{prim_ell},\reg{prim_p})\),
\((\rho,c_H,k)=(\reg{prim_rho},\reg{prim_cH},\reg{prim_k})\) and the strengths
\((r_0,r_1,r_2)=(\reg{prim_r0},\reg{prim_r1},\reg{prim_r2})\). Then \(\tau(r)=(48r-3)/[4(20r-r^2-1)]\), with
\(\tau(r_0)=\reg{tau_r0_frac}\), \(\tau(r_1)=\reg{tau_r1_frac}\) and \(\tau(r_2)=\reg{tau_r2_frac}\). Also
\(\DT(r_0)=1/60\), \(\DT(r_1)=2/3\) and \(\DT(r_2)=169/180\) by \eqref{eq:payoffs}.

\emph{Part (a).} At \(\pi=1/2\), \(M=(1+u)/(2+u)\) and \(\kappa=m=1-M=1/(2+u)\). The left side of (H1) has the
margin \(c_H-B_{r_0}(1/2)=\reg{margin_weakcost_B}\), and the left side of (H2) has the margin
\(k-\DT(r_0)=\reg{margin_kweak_B}\). Neither depends on the slot. The right side of (H1) is \(M>\tau(r_1)\), that is
\(200(1+u)>141(2+u)\), that is \(u>\reg{u_lo_frac}\). Hypothesis (H3) is \(M<\tau(r_2)\), that is
\(5804(1+u)<4245(2+u)\), that is \(u<\reg{u_hi_frac}\). The right side of (H2) is \(\rho(1-M)\DT(r_1)>k\), that is
\((1-M)/6>1/50\), that is \(M<22/25\), that is \(u<\reg{u_kstrong_frac}\). The bound \(u<\reg{u_hi_frac}\) implies it; at
the benchmark slot \(\reg{u_kstrong_frac}-u=\reg{margin_u_kstrong_B}\).
Hence (H1) to (H3) hold if and only if \(\reg{u_lo_frac}<u<\reg{u_hi_frac}\), that is \(\tau(r_1)<M<\tau(r_2)\). At
fixed \(\beta\), \(\lambda=u\beta/(1+u\beta)\) increases strictly in \(u\), so the band is an open interval of
\(\lambda\).

\emph{Part (b).} For \(u<\reg{u_kstrong_frac}\), Step 2 of the proof of Proposition~\ref{prop:ccc} holds at \(r_1\),
and at \(r_2\) because \(\DT(r_2)>\DT(r_1)\). Orders are full at both strengths, and the posteriors are
\((m,1/2,M)\). If \(u<\reg{u_lo_frac}\), then \(M<\tau(r_1)\), no order recruits at \(r_1\), and entry at \(r_1\) is
\(\rho\). If \(\reg{u_hi_frac}\le u<\reg{u_kstrong_frac}\), then \(M\ge\tau(r_2)\), and a purchase recruits at \(r_2\)
under rule C1. For \(u\ge\reg{u_kstrong_frac}\), the right side of (H2) fails at \(r_1\).

At \(u=\reg{u_lo_frac}\), the left side of (H1), both sides of (H2) and (H3) still hold strictly. At \(r_0\), every
correct order earns at most \(\DT(r_0)-k=-1/300<0\), so orders are uniquely zero. At \(r_1\) and \(r_2\), orders are
uniquely full by Step 2, and parts (i) and (iii) hold. A purchase at \(r_1\) leaves the high-cost challenger indifferent, and
rule C1 gives entry \(\Ent_1\). Rule CM allows entry after a purchase anywhere in \([\rho,1]\), so entry at \(r_1\)
can take any value in \([\rho,\Ent_1]\).

\emph{Part (c).} Fix \(1+u=\eu\) and the other benchmark inputs, and vary \(\pi\). The inequality \(M(\pi)>\tau\) is
\(\pi/(1-\pi)>\tau/[(1-\tau)(1+u)]\), that is \(\pi>\tau/[\tau+(1-\tau)\eu]\). The right side of (H1) gives the
lower endpoint of the proposition, and (H3) gives its upper endpoint. The left side of (H1) is
\(B_{r_0}(\pi)<c_H\), that is \(\pi<\tau(r_0)=\reg{tau_r0_frac}\). The left side of (H2) does not depend on \(\pi\).
The right side of (H2) is \(\kappa(\pi)>k/(\rho\DT(r_1))=3/25\): \(m(\pi)>3/25\) is \(\pi>3\eu/(22+3\eu)\), and
\(1-M(\pi)>3/25\) is \(\pi<22/(22+3\eu)\), that is \(\pi\in(\reg{prior_kappa_lo_B},\reg{prior_kappa_hi_B})\)
approximately. This interval contains the other two. The intersection is the interval of the proposition, with
endpoints \(\approx\reg{prior_lo_B}\) and \(\approx\reg{prior_hi_B}\). Without (H3), the upper endpoint is
\(\tau(r_0)\approx\reg{prior_ii_hi}\).
\end{proof}

\paragraph{The region as a result.}
Table~\ref{tab:region} gives the region in decimals, with the margin of the benchmark to each boundary. Two further
prior windows complete it at calibration B. The trade bound \(\kappa(\pi)>3/25\) at \(k=\reg{prim_k}\) holds on
\((\reg{prior_kappa_lo_B},\reg{prior_kappa_hi_B})\). Some trade cost satisfies (H2) when
\(\DT(r_0)<\rho\kappa\DT(r_1)\), that is \(\kappa(\pi)>1/10\), on \((\reg{p2_prior_trade_lo},\reg{p2_prior_trade_hi})\)
(\nolinkurl{P2_auction/derive.py}, run log
\nolinkurl{P2_auction/run.log}). The margins of every hypothesis at
calibration B are in Table~\ref{tab:margins}.
